\anexo{Cronograma de actividades}

El cronograma definitivo deber\'a completarse con las fechas reales de avance, entregas parciales, validaciones y cierre del proyecto. De manera preliminar, se consideran las siguientes etapas:

\begin{enumerate}
    \item Relevamiento del dominio, lectura de bibliograf\'ia y an\'alisis de soluciones existentes.
    \item Definici\'on de requerimientos funcionales y no funcionales.
    \item Dise\~no de arquitectura, modelo de datos y flujos de usuario.
    \item Implementaci\'on del registro de animales, lotes, establecimientos, roles y permisos.
    \item Implementaci\'on de gesti\'on sanitaria, reproductiva, calendario y alertas.
    \item Desarrollo de reportes, indicadores y m\'odulo econ\'omico preliminar.
    \item Integraci\'on de componentes de inteligencia artificial, chatbot y detecci\'on de anomal\'ias.
    \item Pruebas funcionales, validaci\'on con usuarios y ajustes.
    \item Redacci\'on final, revisi\'on bibliogr\'afica y preparaci\'on de la defensa.
\end{enumerate}

%\begin{figure}[ht]
%    \centering
%    \includegraphics[width=0.9\textwidth]{./images/cronograma.png}
%    \caption{Cronograma de actividades.}
%    \label{fig:cronograma}
%\end{figure}
